%% kapa: exam=1 round=30 year=2019 period=2 form=A subject=appraisal_law subject_ko=감정평가관계법규 seq=38 number=38
%% source: Q-Net 감정평가사 기출문제 (cst00309) artlSeq=5205205 file=extracted/1008/5205205/2019년 제30회 감정평가사 1차시험 문제지/2019년 제30회 감정평가사 1차시험 문제지/2019년도 제30회 감정평가사 1차 2교시 A형.hwp
%% conversion: kordoc 4.24.51 -> tools/gen1.py (auto); review=proofread-sonnet; flags=none
\begin{question}{38}
도시 및 주거환경정비법령상 도시ㆍ주거환경정비기본계획(이하 ‘기본계획’)에 관한 설명으로 옳지 않은 것은?
\begin{choices}
  \item 도지사가 대도시가 아닌 시로서 기본계획을 수립할 필요가 없다고 인정하는 시에 대하여는 기본계획을 수립하지 아니할 수 있다.
  \item 정비사업의 계획기간을 단축하는 경우 기본계획의 수립권자는 주민공람과 지방의회의 의견청취 절차를 거쳐야 한다.
  \item 기본계획에는 세입자에 대한 주거안정대책도 포함되어야 한다.
  \item 대도시의 시장이 아닌 시장은 기본계획을 수립하려면 도지사의 승인을 받아야 한다.
  \item 기본계획의 수립권자는 기본계획을 수립하는 경우에 14일 이상 주민에게 공람하여 의견을 들어야 한다.
\end{choices}
\end{question}
\answer{2}
